\documentclass[11pt]{article}
\usepackage[T1]{fontenc}\usepackage{lmodern,amsmath,amssymb,amsthm,booktabs,longtable,array,geometry}
\geometry{margin=.9in}
\newtheorem{theorem}{Theorem}\newtheorem{lemma}{Lemma}
\newcommand{\dd}{\,d}\newcommand{\ellq}{\ell}
\title{Canonical-regime audit: timing, entry geometry and negative weak input}
\author{Independent analytical increment}\date{Version 1, November 2026}
\begin{document}\maketitle
\begin{abstract}
For fixed \(2<\kappa<3\) the nominal strong clock precedes every
target-only timing threshold in the roadmap. Even the zero-noise
right-half tail asymptotic in TAIL is outside its domain at that clock.
We prove a two-sided left-facing angle and mass law for compact interior
Q2 ancestry in the positive-noise comparison. For the leading weak field
we give the receiver contraction with its mass-deficit term and identify
the distinct boundary condition needed for chart return. Neither comparison
supplies initialized canonical entry. The audit separates these conclusions
from the proved \(q^{10}\) chain.
\end{abstract}

\section{Clocks, models and source domains}\label{an04crav1:scope}
Let \(0<q<1,\ S_0=q^\kappa,\ \ell=\log(1/q),\
\lambda=\rho^2,\ \rho\in[13/20,7/10]\), with \(\kappa\) fixed as \(q\downarrow0\).
Denote \(\lambda_- =169/400,\ \lambda_+=49/100\).
Times are normalized. The nominal clock \(T_A=\log(A/S_0)=\kappa\ell+\log A\)
has fixed \(A>0\). It is distinct from a first hitting time in the
trained population. A hat denotes the positive-noise target-only flow.
A bar in the outer comparison denotes its zero-noise reference.

\begin{itemize}
\item F/P give the exact balanced original dynamics for aligned isotropic
initial mass \(S_0>0\). P's early envelopes, including
\(\bar S=S_0e^{(1+2q^2)t}\), compare angles and weights on their
early interval (in particular the alignment guard is \(<\pi\)).
P's separate escape example fixes its parameters and time; it is
not a uniform saturation statement.
\item ARR main uses \(3/5\le\rho\le3/4,\ q\le1/20,\
0<S_0<\delta\le1/100,\ R\ge1,\ qR\le\lambda/2\).
TAIL earlyfraction uses the same small mass cap at
\(T_\delta=\log(\delta/S_0)/(1+2q^2)\). TAIL's exact quadrant
pushforward and its nominal-scale tail corollary refer to a different
model, the zero-noise target-only flow.
\item Q2 cross, transit and outer use a compact fixed interior initial
angle interval in Q2 and compact ratios in \((0,1)\). For transit to
\(qR\) choose \(R>4\rho_+\phi(0)/(1-\rho_+^2)\); fixed Q1 test angles
lie in a compact subset of \((0,\pi/2)\). No training mass guard enters
the scalar target-only clocks. Q2E crossclock covers the closed sector
at positive noise on the ratio box.
\item PROFILE arrival uses fixed sufficiently large \(L\); entry additionally
requires \(T\ge(2/\omega)\ell+C_1,\ q^2T\le1\). The mass law is for
Q1, displacement lower bounds for its upper half, and a finite tail cutoff.
Its joint-power application requires a separately transferred profile and
passage bounds. Q2E's whole-Q2 \(q^3,q^5\) energy equivalences fix
\(T=10\ell,\ S_0=q^{10}\).
\item Q2 originaltracking uses the same initial label, \(q\le1\),
the original input \(a_1\le0\), \(\Delta<\pi\), and full-output defect
envelopes through the claimed time. It gives estimates in terms of those
envelopes without bounding them at saturation.
\item WIT popsystem and P leading describe a given leading characteristic
solution (nonnegative finite mass measures and finite required moments).
The derivative below holds at a frozen state and fixed output coordinate.
WEAK's transit forcing estimate uses \(M=1\), logistic weak mass
\(0<N_0\le N\le n_b<1\), \(Y\ge0\), bounded \(\rho K_s\) and
incoming moments. Its selection theorem also uses nonnegative incoming
input, its moment boundary barrier and strong mean/shape and small-gain
conditions. Neither theorem supplies original-flow entry from Q2.
\end{itemize}

\section{A rational regime map}\label{an04crav1:map}
Write \(a=\rho K(\rho a)>0,\ K(t)=\phi(t)+t\Phi(t)\),
\(P=\Phi(\rho a)\), \(\omega=(1-\lambda)/2,\ d=(1-\lambda P)/2,\
s=d/\omega,\ p=3/s\). PROFILE gives \(1/2<P<153/250\).
Define \(b_-=169/400,\ b_+=49/100,\ f(b)=2/b+4/(1-b)\), and
\[
 L_-=2+169/231,\qquad L_+=250/97+49/51.
\]
\begin{theorem}\label{an04crav1:thresholds}
The following are rational enclosures of the threshold coefficients.
In every row \(2<\kappa<3\) is strictly below the lower endpoint.
The comparison is between logarithmic powers only.
\[
\begin{array}{l|c|l}
\text{event or test}&\text{enclosure}&\text{source and status}\\ \hline
\text{interior Q2 crosses }e_2 &[200/49,800/169]&\text{Q2 cross, target only}\\
\text{last Q1 hits }qL &[1600/231,400/51]&\text{PROFILE arrival, target only}\\
\text{interior Q2 hits fixed Q1 angle}
 &[2/(b_+(1-b_+)),2/(b_-(1-b_-))]&\text{Q2 transit, target only}\\
\text{last Q2 crosses }e_2 &[400/49,1600/169]&\text{Q2E crossclock, target only}\\
\text{interior Q2 hits }qR &[f(b_-),f(b_+)]&\text{Q2 transit, target only}\\ \hline
 E_{\rm loc}>0\text{ threshold}&[2L_-,2L_+]&\text{conditional passage}\\
 pE_{\rm loc}>1\text{ threshold}&[8L_-/3,8L_+/3]&\text{conditional tail charge}
\end{array}
\]
\end{theorem}
\begin{proof}
The first five clocks divide by \(\ell\) and have coefficients
\(2/\lambda,\ 4/(1-\lambda),\ 2/[\lambda(1-\lambda)],\
4/\lambda,\ f(\lambda)\), respectively. A strict comparison with
\(\kappa\) decides the nominal ordering. Leading coefficients alone
do not resolve equality. For all-Q1 arrival PROFILE gives
\(\omega^{-1}\log(1/[q(\cos\alpha+q)])+O(1)\) on the upper half,
including its endpoint. For the last Q2 crossing Q2E gives
\((2/\lambda)\log[(c_0+q)/(q(s_0+q))]+O(1)\),
\(c_0=-\cos\alpha,\ s_0=\sin\alpha\), with the maximum at \(\pi\).
These are different endpoint layers. The formulas are independent of
initial radial scale in the normalized target-only angle equation.
The third coefficient decreases since \(\lambda<1/2\).
The last increases because \(f'>0\) iff \(2\lambda^2>(1-\lambda)^2\),
and \(2(169)^2>231^2\). The other endpoint orders are immediate.

For the last two rows the candidate normalized upper tail is
\(\min(1,(\epsilon/h)^p)\), \(\epsilon=q^{-s}e^{-dT_A}\).
The presumed local shape multiplier has scale
\(G\asymp S_0^{-(1-\lambda)\alpha/\lambda}\),
\(\alpha=\lambda P/2\), from the candidate weak seed
\(N_0\asymp S_0^{1-\lambda}\) and a bounded terminal mass.
This yields \(\epsilon G\asymp q^{E_{\rm loc}}\),
\(E_{\rm loc}=\kappa(1-P)/2-s\). PROFILE's implication requires
the actual entry law, local multiplier, mean, feedback and defect budgets.
The mass charge through an \(O(\ell)\) passage at fixed radius further
requires bounded amplification of the complementary entry mass.
Then \(\mu/q\) costs \(O(q^{pE_{\rm loc}-1}\ell)\).
A shrinking radius \(q^\chi\) and amplification \(q^{-r_m}\) change the
power to \(p(E_{\rm loc}-\chi)-1-r_m\).
Now \(s/(1-P)=1/(1-P)+\lambda/(1-\lambda)\) lies strictly
between \(L_-\) and \(L_+\) by monotonicity. This proves the enclosures
by scalar identities, without assuming independent extremizers exist.
The rational bounds all exceed 3. In particular \(E_{\rm loc}<0\) in
the requested exponent range. A failed upper-bound test supplies no
lower bound on escaped mass or bridge size.
\end{proof}
The fifth threshold solves
\(\gamma_\kappa=1-\lambda(\kappa-4/(1-\lambda))/2=0\);
in the audited range \(\gamma_\kappa>1\).
For orientation within the proved deep chain, \(1600/169<10\).
PROFILE and TAILREP give at ten and \(\chi=1/20\)
\(E_{\rm loc}>3616/6375,\ p(E_{\rm loc}-\chi)-1-r_m\ge4561/35006\)
and \(\Delta_\nu>3/25\) for \(0\le\nu\le1/100\). BULK's secondary
test there has \(\gamma\ge6481/18480,\
\Gamma_\nu\ge289169/924000>3/10\). These are ten-specific margins,
not additional crossing thresholds. Finally \(20^2<5000<20^3\)
gives \(2<\log(5000)/\log20<3\). Only the exponent placement is
asserted at the specified finite pair.

\section{Deep-chain dependence ledger}\label{an04crav1:ledger}
The first failure depends on which statement within an increment is used.
At fixed \(A>0\), TAIL's corollary \texttt{exponent} requires
\[
 (S_0/A)^\omega/q=A^{-\omega}q^{\kappa\omega-1}\to0 .
\]
Instead \(\kappa\omega<3(231/800)=693/800<1\). Its proposed asymptotic
mass is proportional to \(Aq^{3(\kappa\omega-1)}\), unbounded although
the nominal-clock reference mass is bounded by a constant times \(A\)
by TAIL quadrants. Thus the small-tail conclusion itself cannot extend
there. This is an earlier, weaker threshold \(1/\omega\) than the
\emph{all-label} PROFILE arrival threshold \(2/\omega\).
The earlier original-flow ancestral domination is unaffected.
If only initialized all-Q1 profile handoffs are counted, PROFILE's
required time fails and OWPT cannot open its overlap. These distinctions
matter when choosing the S8 starting point.

In the ledger below D means \emph{downstream of the first failure for
the initialized application}. A D row identifies its deep inputs and
where its result is used. Algebraic identities, target-only clocks and
conditional tools retain their stated domains independently. In the deep
rows write \(k=10(1-\lambda),\ h_1=q^{10d-2s},\
\epsilon_h\asymp q^{10d-s}\), \(t_0=t_h=10\ell+O(1)\);
\(H_0\asymp A_q^{-\lambda}q^k\).
Thus the appearance of these symbols in a row includes the dependence
on ten even if the local equation has no explicit exponent.
\small
\begin{longtable}{@{}p{.13\textwidth}p{.79\textwidth}@{}}
\toprule Source & Model, range, and exponent-sensitive contract\\\midrule\endhead
F & Independent. Balanced original dynamics and exact rates at regular times for positive initial mass.\\
P & Independent early envelopes through the alignment guard. \(q^2\log(1/S_0)=\kappa q^2\ell=o(1)\). Model-specific fixed-parameter examples are not steps of this asymptotic implication.\\
ARR & Holds through \(T_\delta\), choose fixed \(\delta\le1/100,\ R=1\). For small \(q\), \(q^\kappa<\delta,\ q\le\lambda/2\). The original incoming mass is at least \(qS_0(S_0/\delta)^{1/16}/3000\); its label set is a backward terminal-sector preimage. No later capture follows.\\
TAIL & Earlyfraction holds: \(q^2\log(\delta/S_0)=o(1)\), and original positive-coordinate ancestry has mass \(\ge e^{-4\delta}\delta\exp[-2q^2\log(\delta/S_0)/(1+2q^2)]/4\). First failure is \texttt{exponent} above, for the zero-noise reference at \(T_A\). Its conditional leading shape and flux identities have no fixed \(\kappa\).\\
Q2 & D. Crossing/transit on compact interior labels are comparison laws for arbitrary \(S_0\); originaltracking is conditional while left-facing. The original weighted defect is unpaid here.\\
PROFILE & D. \(T_A\ge(2/\omega)\ell+C_1\) fails. Its instantaneous outside cost is independent. Tail-inclusive passage needs \(pE_{\rm loc}>1\), and its shrinking-core margin is evaluated at ten.\\
Q2E & D. Endpoint-regularized clocks are comparison laws. The every-label conclusion needs \(4/\lambda<10\); fixed \(T=10\ell\) supplies aggregate \(J\asymp q^3,H\asymp q^5\) and recruited tail.\\
OWPT & D. \(t_m=(2/\omega)\ell+C_m\), \(\sigma=10-2/\omega,\ \zeta=\sigma-1\). The bounds \(S=O(q^\sigma)\), scaled error \(O(q^\zeta)\) need positive powers. Replacing ten by \(2<\kappa<3\) reverses even the mass exponent.\\
SMCP & D. Fixed \(C_m\), \(M_0\asymp e^{C_m}q^\sigma\), shape handoff \(A_q^dq^{10d-s}\). Even at ten this is an obstruction for a uniform linear multiplier on outer labels.\\
MOMC & D. Moving \(C_m=\ell/2\) requires \(\sigma-1/2,\zeta-1/2>0\); passage from this overlap to \(M=1/2\) supplies \(t_h,\epsilon_h\) and the compact core.\\
Q2EP & D. Uses MOMC to bound whole Q2 mass at \(t_h\) by \(C A_q^{-1}q^3\); comparison budgets contain \(q^{10}(1+q^3 e^t)\) and \(q^{10}(e^{\lambda t}+q^5 e^t)\).\\
Q2EC & D. At \(t_h\), bulk has already crossed, latest crossing precedes \(10\ell\) by a diverging margin; \(H_0\), clock ratio \(O(q^{1/8})\), \(c_0\)-threshold \(Kq^{\gamma_{10}}\), and charted recruits.\\
SEP & D. At \(t_h\), secondary mass \(O(q^3)\), scale \(A_q^dh_1\) and index \((1-3\lambda/2)/d\).\\
RET2 & D. Leading stopped mixed-order retention assumes a supplied entry radius, mean budget and passage length. No initialized power law follows from the implication.\\
LOC & D. Leading local clock gain \(N_0^{-\alpha/\lambda}\) requires an entered small core and controlled mean and defects.\\
CLOCK & D. c-form passage requires small integrated forcing, positive amplitude and initial radius times the clock gain small.\\
TAILREP & D. Conditional capped gain and tail moment; takes \(r_m=[10(1-\lambda)(1-P)-2-2\chi]_+\) and the ten-specific positive margins.\\
BULK & D. \(t_0=10\ell+O(1)\), \(H_0\), clock cohort \(c_0\ge q^{1/4}\), \(J_0\le Cq^3\); secondary bound \(q^5+q^{\Gamma_\nu}H\), \(\Gamma_\nu=\gamma_{10}-\nu(k-2)\).\\
RCET & D. Stopped small-ratio transport, comparison dichotomy with regularized endpoints at \(10\ell+O(1)\); retains the bulk clock for vanishing defect.\\
LCRC & D. Joint ledger/closure on \(0\le c\le c_b\) consumes BULK/RCET entry ratios, \(H_0,J_0,h_1\) and retained core/outside budgets.\\
CFR & D. Stopped c-form decomposition uses these cohorts and guards; selected secondary and outside charges initially unpaid.\\
DECG & D. Stopped distance/gain estimates use \(H_0,\epsilon_h,h_1\) and a fixed-chart field allowance.\\
GESG & D. Global envelope and signed-gate estimates, Q3 frozen masses initialized at \(q^{10}\), deep tail scales.\\
S2CP & D. Joint secondary guard for \(q^5+q^{\Gamma_\nu}H\); complement charge \(q^{\Delta_\nu}\ell\) and clock amplification \(O(q^{-k})\).\\
CFA & D. On stopped weak passage uses core cutoff \(q^{\chi+kP/2}\), overshoot and small fixed forcing/strip budgets.\\
Q3SB & D. Pre-entry mass envelope \(q^{10}e^{(1+2q^2)t}\), \(t_m,t_h\), then weak clock with seed \(q^k\).\\
PISB & D. Fixed near-\(\pi\) strip, deep initialization and pre/post half-mass estimates; mass divided by q costs a polylogarithmic factor times \(q^{411/800}\) up to \(c_b\).\\
LSB & D. Late-joiner time shift from deep overlap; comparison mass \(O(q^{10}+q^{12}e^t)\) and original strip gain yielding \(q^{769/4000}\) up to logarithms in mass divided by q.\\
WPT & D. First-exit assembly from \(t_h\), small outside mass and \(H_0\); two-boundary conclusion before RB.\\
RB & D. Exact initialized barrier on \([t_h,t_b]\), with deep entry and Q3 budgets, yields \(H(t_b)\asymp c_b\) and elapsed time \(Q^{-1}[\log(c_b/H_0)+O(1)]\).\\
SI & D. Selection interface inside this passage with \(k=10(1-\lambda)\); boundary field missing in that increment.\\
WBR & D. Weak-side return of the deep clock cohort before \(c_b\); uses the deep chart, positive weak coordinates and cross-field powers.\\
MP & D. Deep clock/logistic matching from \(H=q^{1/8}\); omitted primary amplitude fraction \(O(q^{1/4})\), integrated clock defect \(O(q^{1/8}+q^2\ell)\), phase shift \(t_h-10\ell\), regenerated \(J\le C(q^3+q^2H)\).\\
SRR & D. Strong-reference error on selection interval after the deep return, \(H_0\) and retained radius \(O(q^{1/20})\).\\
WEAK & D for application. Leading \(M=1\) transit and selection have no explicit ten. Sign, boundary and incoming moment conditions must be proved for a new entry.\\
WCS & D. Fixed small weak level \(n\) after deep passage; reached chart \(|u|+|v|\lesssim q^{-603/640}\), incoming moment/phase matching, \(O_n(q^{1/20})\) concentration.\\
CWI & D. Deep fixed-level entry and centered time. Its static leading first-moment estimate is state dependent and independent of \(\kappa\).\\
PSE & D. Original \(q^{10}\) sector-energy control through passage and at fixed weak levels.\\
FWC & D. Every fixed window after \(t_n\) uses WCS/PSE entry, \(M=1+O(q^{11/10})\), logistic mass error \(O(q^{1/8})\) and matched current first moments.\\
RTA & D. Learned rectangle and windows require FWC \(O(q^{1/20})\) distance and full background precision. Initialized reversal/persistence are for ten.\\
BC & D. On these learned windows the pre-saturation output channel has positive order \(q^2\); obstruction uses RTA/FWC original transfer.\\
CPD & D. Same windows and fixed deep partition; producer and exact-lag errors \(O(q^{1/20})\) after scaling.\\
PHS & D. The selected-orbit rational interval is independent of initialization. Original weak-input sign transfer uses CPD on \(J\cap(-\infty,1/1000]\).\\
\bottomrule\end{longtable}\normalsize

\section{Left-facing ancestry at the nominal clock}\label{an04crav1:seed}
\begin{lemma}\label{an04crav1:outer}
Fix \(0<d_-<d_+<\pi/2\), \(\alpha=\pi/2+d_0,\ d_0\in[d_-,d_+]\).
For \(T=\kappa\ell+C\) with bounded \(C\) and \(\kappa>0,\
\kappa\lambda/2<1\), the target-only law is
\[
 \widehat\theta(T)-\pi/2=(1+o(1))\tan(d_0)e^{-\lambda T/2},
 \qquad
 \widehat m(T)=(1+o(1))S_0e^{\lambda T}\sin^2\alpha .
\]
Errors are uniform for the displayed angle compact and parameters
with a fixed positive exponent gap \(1-\kappa\lambda/2\).
In particular the angle is comparable to \(q^{\kappa\lambda/2}\)
through the canonical exponent band on the entire ratio box.
\end{lemma}
\begin{proof}
Use Q2 \texttt{outer}, rather than inferring displacement from a
crossing time alone. In its zero-noise Q2 coordinates
\(\bar w=-\bar U_1=\sqrt{S_0}|\cos\alpha|\) and
\(\bar z=\bar U_2=\sqrt{S_0}\sin\alpha\,e^{\lambda t/2}\).
Take \(L=\ell\). Then
\(\bar c(T)=(\bar w/\bar z)/\sqrt{1+(\bar w/\bar z)^2}
\asymp q^{\kappa\lambda/2}\gg2qL\);
\(\rho\bar s(0)\ge\rho_-\cos d_+>2qL\). The proposition's
\(\epsilon_L=T[3q^2/4+K(-L)/(2L)]\) tends to zero by
\(0<K(-L)\le\phi(L)\). It is below \((\log2)/4\).
Therefore the positive-noise trajectory stays in Q2 through \(T\);
both coordinate logs differ by at most \(\epsilon_L\). Apply
\(\arctan(w/z)\sim w/z\) and \(m=w^2+z^2\). On the canonical band
\(1-\kappa\lambda/2>1-147/200>0\). Neither closed endpoint
belongs to this compact-cohort statement. Q2E's regularized crossing
formula explains why those endpoints cannot be adjoined unconditionally.
\end{proof}
In weak-angle units \(u=(\pi/2-\theta)/q\), this is a negative
coordinate of size \(q^{-1+\kappa\lambda/2}\). The target bulk radial
mass is comparable to \(q^{\kappa(1-\lambda)}\) for a fixed positive
length cohort and bounded clock offset. The orientation and mass law
refer to the comparison only.

\paragraph{Original-flow stop for S8(a).}
For a proposed actual entry time offset \(O(1)\) from \(\kappa\ell\),
Q2 originaltracking compares angular lifts and log weights using
\(f=\mathfrak M X_1e_1+g,\ \mathcal E=\sqrt{a_*}\|g\|_{L^2(P_{\rm mix})}\)
with \(\mathfrak M\ge0,\ a_*=1/2+q^2\). Its zero-initial envelopes solve
\begin{align*}
 \Delta'&=a_*\mathfrak M\Delta+3q\mathfrak M+2\mathcal E,\\
 H_Q'&=a_*H_Q+a_*(1+\mathfrak M)\Delta+2q\mathfrak M+\mathcal E,\\
 L_Q'&=2a_*(\Delta+H_Q)+2q\mathfrak M+2\mathcal E .
\end{align*}
For instance, a sufficient first-exit target up to entry is
\[
 H_Q(t)+\Delta(t)\le\tfrac12\inf_{\alpha}
       (\widehat\theta(t,\alpha)-\pi/2),\qquad
 \Delta<\pi,\qquad L_Q(T)=O(1).
\]
In particular the variation-of-constants integral for \(H_Q(T)\)
has weight \(e^{a_*(T-t)}\) and must be below a fixed fraction of
the terminal \(q^{\kappa\lambda/2}\) gap in this route.
The stated original estimates do not supply this accuracy:
the absolute early bound at fixed \(\delta\) is of order \(\delta\),
whereas \(q^{\kappa\lambda/2}=o(\delta)\).
ARR's cohort is different. A new argument must also identify the
actual strong-entry clock used in this comparison. No original seed law
or passage beyond the first left-facing exit is asserted.

\section{The full weak receiver derivative}\label{an04crav1:slope}
\begin{lemma}\label{an04crav1:contraction}
Fix one state of WIT's leading system with finite nonnegative measures,
finite weak first moments, masses \(M,N\), and \(Y=\int y\,d\mu_s\).
Hold the donors, moments and the receiver output \(v\) fixed.
Write \(F_w^u=u'\) and \(\mathcal W(u)=\int F(u,r)\,d\mu_w(r)\).
For finite \(u\),
\[
 2\partial_u F_w^u=-(u+Y)\phi(u)-(1-M)v\phi(u)
                   -\lambda(1-N)-\mathcal W'(u),\qquad
 \mathcal W'(u)=\phi(u)\int(r-u)_+\,d\mu_w(r)\ge0.
\]
Assume \(Y\ge0\). Put \(b_v=[-(1-M)v]_+\phi(0)/2\).
On the negative half-line and on the indicated subinterval respectively,
\begin{align*}
 \partial_u F_w^u&\le-\tfrac12[\lambda(1-N)-\phi(1)]+b_v
                                      &&(u<0),\\
 \partial_u F_w^u&\le-\tfrac12\lambda(1-N)+b_v
                                      &&(-Y\le u\le0).
\end{align*}
\end{lemma}
\begin{proof}
Differentiate the full equation in WIT \texttt{popsystem}.
P \texttt{overlap} gives \(F_t(t,r)=(r-t)_+\phi(t)\), including the
diagonal by continuity. On compact receiver intervals this is
dominated by \(C(1+|r|)\); differentiation under the integral is
legitimate. Since \(\phi'=-u\phi\), the only adverse contribution
from \(u\phi(u)\) is bounded by \(\phi(1)\), the exact maximum of
\(w\phi(w)\) on \(w\ge0\) at \(w=1\). On \([-Y,0]\),
\((u+Y)\phi(u)\ge0\). These yield the bounds, valid also for non-atomic
weak laws and for a receiver at the boundary \(u=0\).
\end{proof}
For an uncharged estimate require \((1-M)v\ge0\), for example \(M=1\)
or \(M\le1,\ v\ge0\). A sign of \(v\) alone cannot cover \(M>1\).
The reduced rate is strictly positive only if
\(\lambda(1-N)>\phi(1)\). One explicit uniform range is \(0\le N\le1/4\):
the rate is greater than \(\tfrac12[(169/400)(3/4)-1/4]=107/3200\).
The full subinterval rate needs \(N<1\) for strict contraction.
For example, at \(M=1,\ N=1,\ Y=0,\ u=-1\) with weak donors at \(-1\),
the exact slope is \(\phi(1)/2>0\). A uniform positive contraction
on the negative half-line cannot extend through saturation on these
hypotheses. With sign unknown, an integrated charge
\(\int |1-M||v|\phi(0)/2\) bounds the loss in the contraction exponent;
a pointwise rate needs a pointwise charge. For two receivers at the
same \(v\), integrate the derivative along the segment between inputs.
With different output coordinates one must additionally bound the
forcing \(|1-M|\Phi(u)|v_1-v_2|/2\), using the corresponding endpoint
in the split. These are receiver estimates in the same field.

At \(u=0\), \(2F_w^u=\phi(0)-\mathcal W(0)-[Y+(1-M)v]/2\).
A negative slope neither makes this boundary value positive nor proves
arrival at the boundary. To propagate a moving lower face \(-Y\),
one needs the sign of \(F_w^u+Y'\) there. No initialized mass range,
output sign, reached growing chart or inward boundary is inferred from
the lemma.

\section{Open S8 dependencies and the S7 state}\label{an04crav1:handoff}
\begin{enumerate}
\item S8(a) needs a fixed left-facing ancestry at an analytically
matched original entry time, with positive weak-coordinate and radial
mass, input/output angles, imbalance and controlled weights.
Lemma \ref{an04crav1:outer} gives the comparison target.
The weighted full-output defect and survival on the left are open.
\item S8(b) must continue that ancestry through strong learning,
show the negative weak chart is reached with finite-noise errors
controlled on its actual radius, and bound the transient until
entry into the intended weak state. Specify \(M,N,Y,v\) and the
loss of contraction below \(-Y\); prove a crossing/invariance margin
where the consumer requires one. Weak concentration or forcing must
include changes of field, radial weights and output coordinates.
\item S8(c) must give a common clock and actual strong/weak measures:
radial masses and weights, joint input/output profiles on the reached
charts, donor moments, weak seed and its spread, outside-field
cost, and background probe output/rates to the consumer's precision.
In particular the active/inactive split uses the \emph{current} probe
gate. Establish any claimed non-small inactive mass by a lower bound.
\item S7 receives this split state and evolves the active core, inactive
bridge and weak residual producer together, including their backreaction.
Needed outputs include exposure \(A_\zeta=\int(\zeta-x)_+d\mu_s\),
the mixed moment \(\int y(\zeta-x)_+d\mu_s\), compensation and signed
exact lag, and rate-weighted producer contributions on a matched window.
The comparison with coherent depth requires both evolutions. A quantile
description requires established co-order and its preservation in the
same field. Order between different self-consistent populations does
not follow. Errors must resolve the signed lag cancellation and the
\(q^2\)-scale probe effect.
\end{enumerate}
These are open interface obligations. No bridge dynamics, canonical
finite-noise theorem, or new conclusion about OP-W is used.
\end{document}
